\bibliography{custom}

\clearpage
\appendix

\section{Revision Plan for Homework 3}
\label{sec:appendix_revision}

To guide the transition from this Homework~2 manuscript to the final Homework~3 submission, we critically audit our current draft across methodology, empirical scope, and analysis. Table~\ref{tab:revision_plan} details five specific weaknesses identified in our current draft, the exact sections and tables affected, our concrete action plan for Homework~3, and falsifiable completion criteria.

\begin{table*}[t]
\centering
\scriptsize
\renewcommand{\arraystretch}{0.95}
\resizebox{\textwidth}{!}{%
\begin{tabular}{@{}p{1.8cm}p{3.4cm}p{5.5cm}p{4.3cm}@{}}
\toprule
\textbf{Affected Section} & \textbf{Identified Weakness in HW~2 Draft} & \textbf{Concrete Action Plan for Homework~3} & \textbf{Falsifiable Completion Criterion} \\
\midrule
\textbf{W1:} §\ref{subsec:implementation}, §\ref{subsec:main_results}, §\ref{sec:analysis} (Tables~\ref{tab:main_results}--\ref{tab:qualitative}) &
\textbf{Simulated Baseline Numbers Pending Live Colab GPU Execution:} While the complete inference and statistical codebase is implemented and stress-tested, the numerical tables in this draft are populated by the calibrated deterministic simulation backend (\texttt{--backend simulate}) prior to fetching live GPU logs. &
Execute \texttt{src/colab\_runner.ipynb} (\texttt{python -m src.run\_experiment --backend hf --models qwen llama --quantization 4bit --export-latex}) on a Google Colab NVIDIA T4 GPU across all $4{,}200$ prompts, fetch \texttt{results/experiment\_results.json}, and regenerate \texttt{docs/tex/tex/00\_results\_macros.tex} and Table~\ref{tab:qualitative}. &
$100\%$ of entries in \texttt{results/runs\_Qwen\_*.jsonl} and \texttt{results/runs\_meta-llama\_*.jsonl} ($2{,}100$ lines each) originate from live \texttt{--backend hf} T4 GPU execution and are reflected in Tables~\ref{tab:main_results}--\ref{tab:qualitative}. \\
\midrule
\textbf{W2:} §\ref{subsec:rw_generator}, §\ref{subsec:protocol}, §\ref{sec:analysis} &
\textbf{Absence of Conflict-Aware Mitigation Baselines:} The current draft diagnoses severe Truth Override Rates ($\text{TOR} = \qwenPairTwoOverrideAll\%$--\llamaPairTwoOverrideAll\%) in $\text{Pair}_2$, but does not evaluate whether prompt-level or decoding-time interventions can reduce English authority bias. &
Implement and evaluate two inference-time mitigation strategies in \texttt{src/prompts.py} and \texttt{src/models.py}: (1) \emph{Language-Neutral Verification Prompting} (instructing the model to cross-check claims independently of passage language) and (2) \emph{Prior-Anchored Logit Calibration} comparing closed-book vs.\ context-conditioned log-probabilities. &
Add a dedicated mitigation subsection and results table in §\ref{sec:analysis} reporting $\Delta\text{EPR}$ and $\Delta\text{TOR}$ for both interventions across all $600$ counterbalanced pairs per model. \\
\midrule
\textbf{W3:} §\ref{subsec:protocol}, §\ref{subsec:dataset} &
\textbf{Two-Document ($k{=}2$) Pairwise Contradiction Only:} Real-world mRAG pipelines typically retrieve $k \in \{3, 5\}$ passages where a single outdated English passage may conflict with a majority of two or more local-language Korean passages. &
Extend \texttt{src/prompts.py} with a \emph{Majority-Vote Conflict Ablation} ($k{=}3$: $1$ false English passage $d_{\text{en}}(B_i)$ vs.\ $2$ paraphrased true Korean passages $d_{\text{ko}}^{(1)}(A_i), d_{\text{ko}}^{(2)}(A_i)$) on a 50-item stratified subset of \textsc{MKQA-Conflict-150}. &
Report exact Truth Override Rates comparing $1\text{ EN vs.\ }1\text{ KO}$ against $1\text{ EN vs.\ }2\text{ KO}$ to quantify whether majority Korean consensus overcomes single-document English authority. \\
\midrule
\textbf{W4:} §\ref{subsec:controls}, Table~\ref{tab:dataset_examples} &
\textbf{Single-Sentence Template Brevity vs.\ Multi-Sentence Passages:} While single-sentence templates strictly eliminate length and rhetorical confounders, authentic Wikipedia passages embed factual claims inside multi-sentence paragraphs. &
Construct an expanded multi-sentence context condition by adding two neutral, non-conflicting background sentences around $T_{i,\text{en}}$ and $T_{i,\text{ko}}$ and test whether passage length modulates $\text{EPR}$. &
Report paired $\text{EPR}$ and McNemar $\chi^2$ differences between single-sentence and three-sentence passages in §\ref{sec:analysis}. \\
\midrule
\textbf{W5:} §\ref{subsec:dataset}, §\ref{subsec:ablation_category} &
\textbf{Sino-Korean vs.\ Arabic Numeral Realization in Korean:} In \textsc{MKQA-Conflict-150}, $\textsc{Date}$ and $\textsc{Number}$ items use Arabic numerals in both languages, leaving untested how models behave when numbers are spelled out in Hangul (e.g., \ko{사}, \emph{sa} vs.\ \texttt{4}). &
Add a spelled-out Hangul numeral sub-experiment on the 50 $\textsc{Number}$ items comparing Arabic digits against Hangul numeral strings. &
Quantify the exact $\Delta\text{EPR}$ between Arabic-numeral Korean passages and Hangul-spelled Korean passages on the $\textsc{Number}$ category. \\
\bottomrule
\end{tabular}%
}
\vspace{-3pt}
\caption{Structured Revision Plan for Homework~3 detailing five specific weaknesses in the current draft, affected sections, actionable experimental extensions, and falsifiable completion criteria.}
\label{tab:revision_plan}
\vspace{-4pt}
\end{table*}

\section{AI Usage Disclosure}
\label{sec:appendix_ai}

In accordance with the course policy on transparent scientific reporting, we document all uses of generative AI tools across Homework~1 and Homework~2, distinguishing automated assistance from manual author verification:
\begin{itemize}
    \item \textbf{Literature Exploration and Related Work Synthesis (\texttt{NotebookLM}, \texttt{Gemini}):} During Homework~1 and Homework~2, \texttt{NotebookLM} and \texttt{Gemini} were used to summarize the five core reference papers \citep{ki2025,li2025,park2025,qi2025,wang2026} and extract candidate secondary citations from their bibliographies. \emph{Author Verification:} Every citation key, venue, page range, and quantitative claim in §\ref{sec:related_work} and \texttt{docs/tex/custom.bib} was cross-checked directly against the original PDF files in \texttt{assignment/References/}. No unverified or fabricated citations are included.
    \item \textbf{Benchmark Construction and Korean Morphosyntax (\texttt{Antigravity}):} We used the \texttt{Antigravity} coding assistant to write the deterministic filtering script (\texttt{src/build\_dataset.py}) that extracts 150 aligned English--Korean QA pairs from the official Apple MKQA dataset (\texttt{mkqa.jsonl.gz}) and implements the Unicode Hangul coda (\emph{batchim}) and Sino-Korean numeral particle resolver (\texttt{src/korean\_utils.py}). \emph{Author Verification:} We verified $100\%$ of \texttt{mkqa\_example\_id}, \texttt{query\_en}, and \texttt{query\_ko} entries against the official MKQA release and validated all particle resolutions and extraction rules across 1,800 automated assertions in \texttt{src/tests/test\_pipeline.py}.
    \item \textbf{Experimental Pipeline and Manuscript Drafting (\texttt{Antigravity}):} \texttt{Antigravity} assisted in implementing the dual exact-extraction and teacher-forced concept log-probability scoring pipeline (\texttt{src/scoring.py}, \texttt{src/models.py}, \texttt{src/stats.py}), generating the Google Colab runner (\texttt{src/colab\_runner.ipynb}), and formatting the LaTeX tables and TikZ figures in \texttt{docs/tex/tex/}. \emph{Author Verification:} The author reviewed all mathematical derivations (Equations~\ref{eq:pairs}--\ref{eq:mcnemar}), audited the randomly sampled qualitative examples (\texttt{seed=42}), and verified the Revision Plan.
\end{itemize}

\section{Reproducibility, Hyperparameters, and Prompt Templates}
\label{sec:appendix_reproducibility}

\subsection{Codebase and Execution Commands}
\label{subsec:repro_commands}

The complete source code, dataset (\texttt{src/data/mkqa\_conflict\_150.json}), unit test suite (\texttt{src/tests/test\_pipeline.py}), and Google Colab notebook (\texttt{src/colab\_runner.ipynb}) are publicly available at \url{https://github.com/iacantek/nlp-hw}. Table~\ref{tab:hyperparams} lists all experimental hyperparameters and software versions. Every result table and inline macro can be reproduced end-to-end via:
\begin{enumerate}
    \item \textbf{Run 16 Unit Test Suites (1,800+ Assertions):} \texttt{pytest -v}
    \item \textbf{Run Evaluation on Google Colab T4 GPU (or Local Simulation):}
    {\scriptsize
\begin{verbatim}
# Google Colab T4 GPU (Hugging Face 4-bit NF4):
python -m src.run_experiment --backend hf \
  --models qwen llama --quantization 4bit \
  --seed 42 --resume --export-latex
# Fast CPU Validation / Simulation Backend:
python -m src.run_experiment --backend simulate \
  --models qwen llama --seed 42 --export-latex
\end{verbatim}
    }
    \item \textbf{Synchronize LaTeX Macros from Saved JSON Results:}
    {\scriptsize
\begin{verbatim}
python -m src.export_tex \
  --results results/experiment_results.json \
  --output docs/tex/tex/00_results_macros.tex
\end{verbatim}
    }
\end{enumerate}

\begin{table}[t]
\centering
\scriptsize
\renewcommand{\arraystretch}{0.94}
\resizebox{\columnwidth}{!}{%
\begin{tabular}{@{}p{2.5cm}p{4.8cm}@{}}
\toprule
\textbf{Parameter / Component} & \textbf{Specification} \\
\midrule
Repository URL & \url{https://github.com/iacantek/nlp-hw} \\
Benchmark Dataset & \textsc{MKQA-Conflict-150} ($150$ items; $50$ per cat.) \\
Runs per Item / Total & $14$ runs/item $\implies 2{,}100$ runs/model ($4{,}200$ total) \\
Counterbalanced Pairs & $M = 600$ pairs/model ($300$ EN, $300$ KO query) \\
Models Evaluated & \texttt{Qwen/Qwen2.5-7B-Instruct} (7.6B); \texttt{meta-llama/Llama-3.1-8B-Instruct} (8.0B) \\
Target GPU Hardware & Google Colab NVIDIA T4 (16 GB GDDR6 VRAM) \\
Quantization Config & 4-bit NF4 (\texttt{bitsandbytes}), double quant, FP16 \\
Decoding Config & Greedy ($\tau = 0.0$, \texttt{do\_sample=False}), \texttt{max\_new\_tokens=24} \\
Candidate Scoring & Batched teacher-forced length-normalized $\ell_{\text{concept}}$ \\
Random Seed & \texttt{seed = 42} (dataset, simulation, qualitative sample) \\
Python Dependencies & \texttt{torch>=2.1}, \texttt{transformers>=4.43}, \texttt{bitsandbytes>=0.43}, \texttt{scipy} \\
\bottomrule
\end{tabular}%
}
\vspace{-3pt}
\caption{Summary of experimental hyperparameters, hardware specifications, and reproducibility settings.}
\label{tab:hyperparams}
\vspace{-4pt}
\end{table}

\subsection{Verbatim Prompt Templates}
\label{subsec:prompt_templates}

We use parallel concise-answer system and user prompts (\texttt{src/prompts.py}) accompanied below by Revised Romanization transliterations and English glosses:

\paragraph{1. Cross-Lingual \& Monolingual Conflict Prompt ($\ell_q = \text{en}$):}
\begin{quote}
\scriptsize\ttfamily
\textbf{[System]} You are a precise multilingual question-answering assistant. Read the provided documents carefully and answer the question using ONLY the concise entity, date, or number. Do not include explanations or punctuation.\\[2pt]
\textbf{[User]} Document 1: \{doc1\} \quad Document 2: \{doc2\}\\
Question: \{query\_en\}\\
Answer (concise entity, date, or number only):
\end{quote}

\paragraph{2. Cross-Lingual \& Monolingual Conflict Prompt ($\ell_q = \text{ko}$):}
\begin{quote}
\scriptsize\ttfamily
\textbf{[System]} \ko{당신은 정확한 다국어 질의응답 도우미입니다. 제공된 문서를 주의 깊게 읽고, 설명이나 문장 부호 없이 핵심 개체명, 연도 또는 숫자만으로 간결하게 질문에 답하세요.}\\
\normalfont\scriptsize\emph{(Dangsin-eun jeonghwakhan dagugeo jirui-eungdap doumi-imnida. Jegongdoen munseo-reul juui gipge ikgo, seolmyeong-ina munjang buho eopsi haeksim gaechemyeong, yeondo ttoneun sutja-man-euro gangyeolhage jilmun-e daphaseyo. / ``You are an accurate multilingual QA assistant. Read the provided documents carefully and answer concisely with only the core entity, year, or number.'')}\\[2pt]
\ttfamily
\textbf{[User]} \ko{문서} 1: \{doc1\} \quad \ko{문서} 2: \{doc2\}\\
\ko{질문}: \{query\_ko\} \quad \ko{정답 (핵심 개체명, 연도 또는 숫자만 답하세요)}:\\
\normalfont\scriptsize\emph{(Munseo 1/2; Jilmun; Jeongdap: ``Document 1/2; Question; Answer:'')}
\end{quote}

\paragraph{3. Closed-Book Parametric Prior Prompt ($\ell_q \in \{\text{en}, \text{ko}\}$):}
Omits \texttt{Document 1} and \texttt{Document 2} and instructs the model to answer directly from parametric knowledge in the query language using only the concise entity, date, or number.
